\clearpage
\fancyhead[L]{\small CERES.}
\begin{center}
CERES.

\rule{0.6\linewidth}{0.4pt}

Monatliche Correspondenz. Vol. XI. pp. 283–292. March 1805.

\rule{0.6\linewidth}{0.4pt}
\end{center}

{\rightskip=0pt plus 1em
The period when Ceres was visible is now over as well. Even with larger equatorial instruments, observing it would likely be difficult, since it would have to be observed a considerable distance from the meridian. We must now give up following this wandering star across the sky for six to seven months and aim only to locate it again as soon as it reappears. The elements that Dr. \textsc{Gauss} has improved once more and the ephemeris calculated from them will make this possible without difficulty, since the orbit of Ceres can now be taken as very accurately determined. We give here everything else that may make it easier to observe in future.\par}

The ephemeris of the geocentric path of Ceres for September and October, calculated from the ninth set of elements, gave its right ascension nine minutes too small and its declination four minutes too large. Dr. \textsc{Gauss} therefore made use of the three oppositions of Ceres in 1802, 1803 and 1804 and based the following tenth set of elements on them.
\SourceLeaf{267}
\begin{center}
\begingroup\fontsize{9}{11}\selectfont\setlength{\tabcolsep}{3pt}\renewcommand{\arraystretch}{1.1}
\begin{tabular}{lr}
Epoch 1804, Seeberg meridian . . . . . . . . & $312^\circ\ 1'\ 33''5$ \\
Daily motion . . . . . . . . . . . . . & $771''0524$ \\
Aphelion . . . . . . . . . . . . . . . & $326^\circ\ 26'\ 3''1$ \\
Eccentricity . . . . . . . . . . . . . . . & $0.0784757$ \\
Logarithm of the semiaxis . . . . . . . . . & $0.04420004$ \\
$\AscendingNode$ 1804 . . . . . . . . . . . . . . . . . & $80^\circ\ 59'\ 12''$ \\
Inclination . . . . . . . . . . . . . . . . . & $10\ 37\ 45$ \\
\end{tabular}\endgroup
\end{center}

\indent From these elements Dr. \textsc{Gauss} calculated the geocentric path of this planet from 28 July 1805 to 24 May 1806 [and gives the right ascension and declination to the minute, the distance from Earth to three decimal places and the brightness to five decimal places, for 12 o’clock Seeberg time at three-day intervals.]

—\quad—\quad—

He [\textsc{Gauss}] wrote to us that \textsc{Bessel}, whose paper on the orbit of the comet of 1607 (M. C. 1804, vol. X, p. 425 f.) had earned him such distinction, helped him calculate this ephemeris by supplying all the necessary positions of the Sun. Surely anyone who knows what calculations are needed to determine the elements of a planet and then to derive each position from them must marvel that one man could complete so many laborious calculations in so little time. —\quad—\quad—

\hfill[\textsc{von Zach.}]
